\section{Discussão}

A diferença entre estágios no mesmo eletrodo foi a representação mais informativa: usar apenas o
estágio final, ou só descritores resumidos do pico, reduziu o desempenho. Isso é coerente com o
princípio do sensor, em que o sinal de interesse é a mudança provocada pela hibridização, e não o
valor absoluto da corrente, que varia entre eletrodos e lotes. Na prática, o protocolo de
aquisição precisa garantir que os três estágios sejam medidos no mesmo eletrodo e registrados na
planilha de metadados; réplicas incompletas não podem ser classificadas.

O controle de qualidade automático se mostrou necessário já nos primeiros arquivos reais do
projeto, em que medições com sobrecarga do canal de corrente e com o potencial medido muito
diferente do aplicado produziam curvas sem significado eletroquímico. Sem esse filtro, essas
medições entrariam no treino como se fossem amostras. \preencher{quantificar com os dados reais}

A validação agrupada por amostra e a separação do teste antes de qualquer escolha evitam o
vazamento de informação entre treino e teste, problema frequente em aplicações de aprendizado de
máquina \citep{kapoor2023}. Como dezenas de combinações foram comparadas, a acurácia da melhor
combinação na validação cruzada tende a ser otimista; o valor de referência é o do conjunto de teste.

\paragraph{Limitações.} \preencher{número de empresas, tipos de superfície e bactérias; variação
entre lotes de eletrodos; necessidade de validação em campo (objetivo 7 do projeto); comparação com
métodos rápidos comerciais, se houver}
